\subsubsection{\texttt{aptknt}}

\texttt{aptknt} 根据给出的数据位点以及样条阶数生成合适的B样条基函数节点序列

用法为\texttt{knots = aptknt( tau, k )} 或\texttt{[knots, k] = aptknt( tau, k )}

\begin{itemize}
    \item \texttt{tau}为一个单增的序列，并满足$\forall i,\ tau(i) < tau(i+k-1)$
    \item k为生成的样条阶数，$k \geq 2$
    \item knots 为生成的节点序列
    \item 返回值k为可选项，为实际能够实现的样条阶数，因为所给数据位点不一定能够满足实现k阶样条的条件
\end{itemize}

生成算法为：首尾节点重复k次，中间节点采用$tau^*(i) = \frac{\sum_{j=i+1}^{i+k-1}}{k-1}tau(j)$计算

即\texttt{aptknt( tau, k )}等价于\texttt{augknt( [tau(1),aveknt(tau,k),tau(end)],k )}

样例\textbf{I.}

输入：
\begin{lstlisting}[language = matlab]
    knots = aptknt( [1 2 3 4 5 6], 3 )
\end{lstlisting}

 输出为：
\begin{lstlisting}[language = matlab ]
    knots = 

        1.0000    1.0000    1.0000    2.5000    3.5000    4.5000    6.0000    6.0000    6.0000
\end{lstlisting}

样例\textbf{II.}

输入：
\begin{lstlisting}[language = matlab]
    [knots, k] = aptknt( [1 2 3 4 5 6], 7 )
\end{lstlisting}

 输出为：
\begin{lstlisting}[language = matlab ]
    knots = 

        1   1   1   1   1   1   1   6   6   6   6   6   6   6

    k = 

        1x1 int32

        6
\end{lstlisting}
